% ============================================================
%  КУРСОВА РОБОТА — XeLaTeX
% ============================================================

\documentclass[14pt,a4paper]{extreport}

\usepackage{fontspec}
\setmainfont{Times New Roman}

\usepackage{polyglossia}
\setdefaultlanguage{ukrainian}

\usepackage[a4paper,top=2.5cm,bottom=2.5cm,left=3.5cm,right=1.5cm]{geometry}

\usepackage{setspace}
\onehalfspacing

\usepackage{indentfirst}
\setlength{\parindent}{1.25cm}

\usepackage{titlesec}

\titleformat{\chapter}[block]{\bfseries\centering\MakeUppercase}{}{0pt}{}
\titleformat{\section}[block]{\bfseries}{\thesection}{0.5em}{}

\renewcommand{\thesection}{\arabic{section}}
\renewcommand{\thesubsection}{\thesection.\arabic{subsection}}

\usepackage{fancyhdr}
\pagestyle{fancy}
\fancyhf{}
\fancyhead[R]{\thepage}

% команды как в шаблоне
\newcommand{\myintro}[1]{
\begin{center}
\textbf{\MakeUppercase{#1}}
\end{center}
\vspace{10pt}
}

\newcommand{\mysection}[3]{
\section*{#2. #3}
\addcontentsline{toc}{section}{#2. #3}
}

\begin{document}

% ============================================================
%  ТИТУЛЬНА СТОРІНКА
% ============================================================

\thispagestyle{empty}

\begin{center}
Бердичівський педагогічний фаховий коледж Житомирської обласної ради
\end{center}

\vspace{10pt}

\begin{center}
Циклова комісія дошкільної освіти
\end{center}

\vspace{20pt}

\begin{center}
\textbf{КУРСОВА РОБОТА}
\end{center}

\begin{center}
з методики образотворчої діяльності
\end{center}

\vspace{10pt}

\begin{center}
\textbf{на тему:}

\textbf{«Методика навчання зображення транспорту  
в старшому дошкільному віці»}
\end{center}

\vspace{40pt}

Студента III курсу групи Д-31  

галузі знань: 01 Освіта / Педагогіка  

спеціальності: 012 Дошкільна освіта  

ПІБ СТУДЕНТА

\vspace{15pt}

Керівник: ПІБ КЕРІВНИКА  

викладач педагогіки

\vfill

\begin{center}
Бердичів — 2026
\end{center}

\newpage

% ============================================================
%  ЗМІСТ
% ============================================================

\tableofcontents

\newpage

% ============================================================
%  ВСТУП
% ============================================================

\myintro{ВСТУП}

Дошкільний вік є важливим етапом розвитку особистості дитини. Саме в цей період активно формуються пізнавальні процеси, розвиваються творчі здібності, удосконалюється дрібна моторика та зростає інтерес до пізнання навколишнього світу. Важливу роль у цьому процесі відіграє образотворча діяльність.

Образотворча діяльність допомагає дитині передавати власні враження про навколишній світ, розвиває уяву, мислення та емоційну сферу. Під час малювання діти знайомляться з формою предметів, їхніми пропорціями та кольором.

Старший дошкільний вік є перехідним етапом між дошкільним дитинством і навчанням у школі. У цей період формуються важливі психічні новоутворення, розвивається довільна поведінка та інтерес до навчальної діяльності.

Одним із важливих об'єктів зображення для дітей є транспорт. Діти щодня спостерігають автомобілі, автобуси, потяги та інші транспортні засоби. Це викликає у них інтерес і бажання відтворити побачене у малюнках.

Актуальність теми дослідження зумовлена необхідністю розвитку творчих здібностей дітей дошкільного віку та удосконалення методики навчання малювання.

Мета роботи — дослідити методику навчання зображення транспорту в старшому дошкільному віці.

Завдання дослідження:

1. Проаналізувати психолого-педагогічну літературу з проблеми дослідження.
2. Охарактеризувати особливості розвитку дітей старшого дошкільного віку.
3. Розглянути значення образотворчої діяльності для розвитку дошкільників.
4. Проаналізувати методику навчання дітей зображення транспорту.

Об’єкт дослідження — образотворча діяльність дітей дошкільного віку.

Предмет дослідження — методика навчання зображення транспорту.

\newpage

% ============================================================
%  РОЗДІЛ 1
% ============================================================

\chapter*{РОЗДІЛ 1 \\ ТЕОРЕТИЧНІ ОСНОВИ ОБРАЗОТВОРЧОЇ ДІЯЛЬНОСТІ ДОШКІЛЬНИКІВ}

\addcontentsline{toc}{chapter}{РОЗДІЛ 1}

\mysection{}{1}{Психологічні особливості дітей старшого дошкільного віку}

Старший дошкільний вік характеризується активним розвитком психічних процесів дитини. У цей період розвиваються мислення, пам'ять, увага та уява.

Пам’ять відіграє важливу роль у пізнавальній діяльності людини. Вона забезпечує закріплення, збереження та відтворення попереднього досвіду.

У старшому дошкільному віці зростає роль довільної поведінки. Діти починають краще контролювати свої дії та виконувати поставлені завдання.

\mysection{}{1}{Значення образотворчої діяльності}

Образотворча діяльність є важливим засобом розвитку творчих здібностей дітей. У процесі малювання розвивається уява, мислення та емоційна сфера.

Діти вчаться спостерігати за предметами навколишнього світу, аналізувати їх форму та передавати її у малюнку.

Крім того, малювання сприяє розвитку дрібної моторики рук, що є важливим для подальшого навчання письма.

\newpage

% ============================================================
%  РОЗДІЛ 2
% ============================================================

\chapter*{РОЗДІЛ 2 \\ МЕТОДИКА НАВЧАННЯ ЗОБРАЖЕННЯ ТРАНСПОРТУ}

\addcontentsline{toc}{chapter}{РОЗДІЛ 2}

\mysection{}{2}{Особливості навчання дітей малювання транспорту}

Транспорт є складним об'єктом для зображення, оскільки складається з кількох частин.

Тому навчання дітей повинно відбуватися поступово.

На першому етапі діти знайомляться з різними видами транспорту та їхніми основними частинами.

Педагог може використовувати ілюстрації, іграшкові моделі транспорту або спостереження за реальними транспортними засобами.

\mysection{}{2}{Методи навчання}

Під час навчання використовуються різні методи:

Пояснення та показ — педагог демонструє послідовність створення малюнка.

Спостереження — діти розглядають транспортні засоби та їх зображення.

Поетапне малювання — спочатку зображується основна форма, потім додаються деталі.

Творчі завдання — діти створюють власні малюнки транспорту.

\newpage

% ============================================================
%  ВИСНОВКИ
% ============================================================

\myintro{ВИСНОВКИ}

Образотворча діяльність є важливим засобом розвитку дітей дошкільного віку.

У процесі малювання розвиваються мислення, уява та сприймання.

Навчання зображення транспорту сприяє розвитку творчих здібностей дітей, їх спостережливості та просторового мислення.

Правильно організована педагогічна діяльність сприяє ефективному формуванню образотворчих навичок у дітей.

\newpage

% ============================================================
%  СПИСОК ЛІТЕРАТУРИ
% ============================================================

\myintro{СПИСОК ВИКОРИСТАНИХ ДЖЕРЕЛ}

\begin{enumerate}

\item Варій М. Й. Загальна психологія.

\item Дуткевич Т. В. Дошкільна психологія.

\item Павелків Р. В., Цигипало О. П. Дитяча психологія.

\item Палій А. А. Методи діагностики психічного розвитку дітей.

\item Марінушкіна О. Є. Порадник практичного психолога.

\item Діагностика готовності дітей до школи.

\end{enumerate}

\end{document}